\selectlanguage{english}
\begin{abstract}
The Square Kilometer Array (SKA) will be the first observatory to produce three-dimensional maps of the $21~\si{cm}$ line of neutral hydrogen during Cosmic Dawn and the Epoch of Reionization. Extracting information from its incredibly complex signal will require modern techniques of machine learning. This work applies simulation-based inference to train neural networks for the analysis of future SKA data. The method employs a convolutional neural network (CNN) to summarize the data to highly informative summaries and infers parameters of interest from the summary using an invertible normalizing flow. The same flow model is used to assess the informativeness of the CNN summary with an estimate of the mutual information with the parameter distribution. This thesis analyzes the robustness of this estimate during inference of the neutral hydrogen fraction and of six cosmological parameters used in the simulations.
\end{abstract}

\selectlanguage{german}
\begin{abstract}
Das Square Kilometer Array (SKA) wird das erste Observatorium sein, welches dreidimensionale Daten zur $21~\si{cm}$-Linie des neutralen Wasserstoffs während der Epochen der Kosmischen Renaissance und der Reionierung erstellen kann. Das Extrahieren von Information aus diesem äußerst komplexen Signal erfordert moderne Methoden des maschinellen Lernens. Diese Arbeit verwendet simulationsgestützte Inferenz, um neuronale Netze für die Analyse zukünftiger SKA-Daten zu trainieren. Dabei wird ein Convolutional Neural Network (CNN) eingesetzt, um die Daten in hochinformative Zusammenfassungen zu komprimieren. Aus diesen werden mithilfe eines Invertible Normalizing Flow (INF) kosmologische Parameter inferiert. Dasselbe INF wird verwendet, um die Informativität der CNN-Zusammenfassung anhand einer Abschätzung der Mutual Information mit der Parameterverteilung zu bewerten. Diese Arbeit analysiert die Robustheit dieser Schätzung für die Inferenz des neutralen Wasserstoffanteils sowie von sechs in den Simulationen verwendeten kosmologischen Parametern.
\end{abstract}